% from ega4-17-fr.tex
\oldpage[IV-4]{97}

\begin{enumerate}
  \item[a)] $T$ est lisse sur $S$ au point $t$, de dimension relative $d+\sum_{i\in I}c_i$.
  \item[b)] L'homomorphisme
  \[
  \rho\otimes1:\sh{J}\otimes_{\sh{O}_X}k(t)\to\bigoplus_{i\in I}(\Omega^1_{Y_i/S}\otimes_{\sh{O}_{Y_i}}k(t))
  \]
  est injectif.
\end{enumerate}
Lorsque $t$ est rationnel sur le corps $k(s)$ (où $s$ est l'image de $t$ dans $S$), ces conditions sont aussi équivalentes à la suivante:
\begin{enumerate}
  \item[b')] L'homomorphisme transposé de $\rho\otimes1$
  \[
  \bigoplus_{i\in I}\mathrm{T}_{Y_i/X}(y_i)\to(\mathrm{T}_{X/S}(x))^I/\delta(\mathrm{T}_{X/S}(x))
  \]
  (où $\delta$ est l'application diagonale) est surjectif.
\end{enumerate}
Il suffit de remarquer que $X^I$ est lisse sur $S$ et de dimension relative $d\cdot\operatorname{Card}(I)$, et $Y$ lisse sur $S$ de dimension relative $d\cdot\operatorname{Card}(I)+\sum_{i\in I}c_i$.

Lorsque les conditions équivalentes de (17.13.10) sont vérifiées, on dit encore que les $f_i$ forment une \emph{famille de morphismes transversaux au point $t$, relativement} à $S$. On voit encore que l'ensemble des points $t\in T$ où cela a lieu est \emph{ouvert} dans $T$. La remarque (17.13.4, (iii)) montre alors que si $X$ et les $Y_i$ sont plats et localement de présentation finie sur $S$, et de plus si $T$ est \emph{propre} sur $S$, l'ensemble des $s\in S$ tels que les $f_i$ forment une famille de morphismes transversaux en \emph{tous} les points de $T_s$, relativement à $S$, est \emph{ouvert} dans $S$.
\end{corollary}

\begin{env}[17.13.11]
\EGAReaderTarget{ega4.c04.statement.17-13-11}{IV.17.13.11}
Considérons en particulier le cas, généralisant (17.13.7), où les $Y_i$ ($i\in I$) sont des \emph{sous-préschémas} de $X$, les $f_i$ étant les injections canoniques, de sorte que $T=\inf_{i\in I}(Y_i)$ est encore le sous-préschéma « intersection » des $Y_i$, et $t=y_i=x$ pour tout $i$; au lieu de dire que les $f_i$ forment une famille de morphismes transversaux au point $x$, on dit encore que \emph{les $Y_i$ se coupent transversalement au point $x$} (relativement à $S$). La condition a) de (17.13.10) s'exprime encore en la relation qui généralise (17.13.7.1)
\[
\label{IV.17.13.11.1}
\tag{17.13.11.1}
\operatorname{codim}_x(T_s,X_s)=\sum_{i\in I}\operatorname{codim}_x((Y_i)_s,X_s).
\]
En outre, on a la propriété suivante, qui étend (17.13.8.1), et en donne une autre démonstration quand $I$ a 2 éléments:
\end{env}

\begin{corollary}[17.13.12]
\label{IV.17.13.12}
Lorsque les $Y_i$ se coupent transversalement au point $x$, on a un isomorphisme canonique
\[
\label{IV.17.13.12.1}
\tag{17.13.12.1}
\sh{N}_{T/X}\xrightarrow{\sim}\bigoplus_{i\in I}(\sh{N}_{Y_i/X}\otimes_{\sh{O}_{Y_i}}\sh{O}_T).
\]
\end{corollary}

On peut se borner au cas où les $Y_i$ sont des sous-préschémas fermés de $X$, définis par des Idéaux quasi-cohérents $\sh{J}_i$, de sorte que $T$ est défini par l'Idéal $\sh{J}=\sum_i\sh{J}_i$. Par définition du faisceau conormal à une immersion (16.1.3), l'homomorphisme canonique $\bigoplus_i\sh{J}_i\to\sh{J}$ donne, par passage aux quotients, un homomorphisme \emph{surjectif}
\[
\label{IV.17.13.12.2}
\tag{17.13.12.2}
\bigoplus_{i\in I}(\sh{N}_{Y_i/X}\otimes_{\sh{O}_{Y_i}}\sh{O}_T)\to\sh{N}_{T/X}.
\]

